\chapter{Implementation and Testing}
\label{ch:implementation-testing}

This chapter describes the current implementation status of the CoWriteIA system and the testing strategies applied to validate its functionality. The focus of this phase was to implement the core backend services, major APIs, and conduct rigorous black box testing to ensure that all user-facing functionalities behave correctly according to the system requirements.

The implementation described in this chapter reflects the progress achieved up to the current iteration of the project.

\section{Algorithm Design}
\label{sec:algorithm-design}

The system implements a Retrieval-Augmented Generation (RAG) pipeline to enable intelligent document-based responses. The algorithm focuses only on the core functional requirements: document processing, semantic indexing, context retrieval, and AI response generation.

The CoWriteIA implementation in this repository realizes a semantic indexing + RAG (retrieval-augmented generation) pipeline. The implemented components (and their primary locations in the codebase) are:

\begin{itemize}
    \item Document chunking and preprocessing
    \begin{itemize}
        \item \texttt{backend/app/services/text\_extraction\_service.py :: TextExtractionService.chunk\_text / \_create\_text\_chunk}
    \end{itemize}
    \item Embedding generation and vector storage
    \begin{itemize}
        \item \texttt{backend/app/services/embedding\_service.py :: EmbeddingService.generate\_single\_embedding, generate\_embeddings, store\_text\_chunks\_with\_embeddings}
        \item \texttt{backend/app/services/chroma\_service.py :: ChromaService.add\_embeddings, search\_similar}
        \item \texttt{backend/app/repositories/text\_chunk\_repository.py :: update\_embedding, get\_chunks\_without\_embeddings, vector\_search / vector\_search\_chroma}
    \end{itemize}
    \item Context retrieval (semantic similarity)
    \begin{itemize}
        \item \texttt{backend/app/services/rag\_context\_service.py :: RAGContextService.assemble\_context}
        \item \texttt{backend/app/services/search\_service.py :: SearchService.semantic\_search, find\_similar\_content}
    \end{itemize}
    \item AI response generation (RAG orchestration and assistant)
    \begin{itemize}
        \item \texttt{backend/app/services/ai\_chat\_service.py :: AIChatService.chat}
        \item \texttt{backend/app/services/copilot\_service.py :: CopilotService.generate\_suggestion}
    \end{itemize}
\end{itemize}

The following subsections contain concise algorithm descriptions derived from the actual implementation and precise pseudocode that maps to the repository functions.

\subsection{High-level Flow}
\begin{enumerate}
    \item User uploads file via endpoints in \texttt{backend/app/api/v1/endpoints/files.py}.
    \item File processing triggers text extraction and chunking:
    \begin{itemize}
        \item \texttt{TextExtractionService.extract\_text\_from\_file(...)} and \texttt{TextExtractionService.chunk\_text(...)}
        \item Chunks are created as \texttt{TextChunk} models and persisted via \texttt{TextChunkRepository}.
    \end{itemize}
    \item Embeddings are generated for chunks:
    \begin{itemize}
        \item \texttt{EmbeddingService.generate\_embeddings(...)} / \texttt{generate\_single\_embedding(...)}
        \item Embeddings saved both in chunk metadata (MongoDB) and in ChromaDB via \texttt{ChromaService} when available.
    \end{itemize}
    \item Semantic search / RAG:
    \begin{itemize}
        \item \texttt{RAGContextService.assemble\_context(...)} calls embedding/search services to retrieve top-$k$ chunks.
        \item \texttt{SearchService} handles hybrid/fallback logic (ChromaDB preferred, fallback to MongoDB vector search).
    \end{itemize}
    \item LLM generation:
    \begin{itemize}
        \item \texttt{AIChatService.chat(...)} or \texttt{CopilotService.generate\_suggestion(...)} build a prompt from the retrieved context and call the configured LLM (\texttt{GroqService} or other).
    \end{itemize}
    \item Results returned to the client; metadata and analytics stored as needed.
\end{enumerate}

\subsection*{Example short sequence (mapping code $\rightarrow$ runtime)}
\begin{enumerate}
    \item User uploads file $\rightarrow$ API endpoint in \texttt{backend/app/api/v1/endpoints/files.py} triggers \texttt{process\_uploaded\_file}.
    \item \texttt{process\_uploaded\_file} calls \texttt{TextExtractionService.chunk\_text} to create chunks.
    \item Indexing background task (or immediate synchronous call) calls \texttt{EmbeddingService.generate\_embeddings} for chunks.
    \item \texttt{EmbeddingService} stores embeddings via \texttt{TextChunkRepository} and \texttt{ChromaService}.
    \item Later, user sends search/query $\rightarrow$ \texttt{RAGContextService.assemble\_context} retrieves top chunks using \texttt{ChromaService} or \texttt{SearchService}.
    \item \texttt{AIChatService.chat} builds prompt with context and calls \texttt{GroqService} (LLM); response persisted and returned.
\end{enumerate}

\subsection{Pseudocode: Document Chunking and Preprocessing}

The system splits extracted text into sentence-aware chunks within a target size and optionally applies overlap. Each chunk is materialized using the implementation in \texttt{backend/app/services/text\_extraction\_service.py}.

\begin{algorithm}[H]
\caption{ChunkTextAndCreateChunks}
\begin{algorithmic}[1]
\Require \textit{file\_id}, \textit{project\_id}, \textit{full\_text}, \textit{max\_chunk\_size}, \textit{overlap} (optional), \textit{start\_chunk\_index} (default 0)
\Ensure List of \texttt{TextChunk} models \textit{chunks}
\State sentences $\leftarrow$ \texttt{TextExtractionService.\_split\_into\_sentences}(full\_text)
\State chunks $\leftarrow$ [\,]
\State current\_chunk\_sentences $\leftarrow$ [\,];\; current\_word\_count $\leftarrow$ 0;\; chunk\_index $\leftarrow$ start\_chunk\_index
\For{each sentence $s$ in sentences}
    \State sentence\_words $\leftarrow$ \texttt{count\_words}($s$)
    \If{current\_word\_count + sentence\_words $>$ max\_chunk\_size \textbf{and} current\_chunk\_sentences not empty}
        \State chunk\_text $\leftarrow$ \texttt{join}(current\_chunk\_sentences, " ")
        \State start\_pos $\leftarrow$ compute absolute start in \textit{full\_text}
        \State chunk\_model $\leftarrow$ \texttt{TextExtractionService.\_create\_text\_chunk}(file\_id, project\_id, chunk\_text, chunk\_index, start\_pos, current\_word\_count)
        \State append chunk\_model to chunks
        \State chunk\_index $\leftarrow$ chunk\_index + 1
        \If{overlap configured}
            \State current\_chunk\_sentences $\leftarrow$ take last \textit{overlap} words/sentences of current\_chunk\_sentences
            \State current\_word\_count $\leftarrow$ \texttt{count\_words}(current\_chunk\_sentences)
        \Else
            \State current\_chunk\_sentences $\leftarrow$ [\,];\; current\_word\_count $\leftarrow$ 0
        \EndIf
    \EndIf
    \State append $s$ to current\_chunk\_sentences;\; current\_word\_count $\leftarrow$ current\_word\_count + sentence\_words
\EndFor
\If{current\_chunk\_sentences not empty}
    \State chunk\_text $\leftarrow$ \texttt{join}(current\_chunk\_sentences, " ")
    \State start\_pos $\leftarrow$ compute absolute start in \textit{full\_text}
    \State chunk\_model $\leftarrow$ \texttt{TextExtractionService.\_create\_text\_chunk}(file\_id, project\_id, chunk\_text, chunk\_index, start\_pos, current\_word\_count)
    \State append chunk\_model to chunks
\EndIf
\State \Return chunks
\end{algorithmic}
\end{algorithm}
\FloatBarrier

\textbf{Notes:}
\begin{itemize}
    \item Uses sentence-aware splitting and target size with optional overlap per \texttt{TextExtractionService}.
    \item Returns \texttt{TextChunk} Pydantic models; persistence occurs later (e.g., in \texttt{store\_text\_chunks\_with\_embeddings}).
\end{itemize}

\subsection{Pseudocode: Embedding Generation and Storage}

This stage encodes chunked text into dense vectors and persists both the metadata and vectors. It maps to \texttt{EmbeddingService} and \texttt{ChromaService} usage in the codebase.

\begin{algorithm}[H]
\caption{GenerateAndPersistEmbeddings}
\begin{algorithmic}[1]
\Require \textit{project\_id}, \textit{file\_id}, \textit{list\_of\_chunks} (content strings)
\Ensure List of created/updated \texttt{TextChunk} records with embeddings
\State chunk\_texts $\leftarrow$ [\,\texttt{chunk.content} for \textit{chunk} in \textit{list\_of\_chunks}\,]
\State embeddings $\leftarrow$ \texttt{EmbeddingService.generate\_embeddings}(chunk\_texts)
\Comment{Runs threadpool synchronous encoding; model loads lazily}
\State created\_chunks $\leftarrow$ [\,]
\For{each $(chunk\_info, embedding)$ in \texttt{zip}(list\_of\_chunks, embeddings)}
    \State chunk\_model $\leftarrow$ \texttt{build TextChunk}(\,
        file\_id $=$ file\_id,\; project\_id $=$ project\_id,\;\
        content $=$ chunk\_info.content,\; start\_position $= \dots$,\; end\_position $= \dots$,\;\
        chunk\_index $=$ chunk\_info.chunk\_index,\; word\_count $= \dots$,\;\
        embedding\_vector $=$ embedding\,)
    \State saved\_chunk $\leftarrow$ \texttt{TextChunkRepository.create}(chunk\_model)
    \Comment{Persist metadata in MongoDB}
    \State append saved\_chunk to created\_chunks
\EndFor
\If{\texttt{ChromaService} available}
    \State chroma\_entries $\leftarrow$ map \textit{created\_chunks} to \{\texttt{id}, \texttt{embedding}, \texttt{metadata}\}
    \State \texttt{ChromaService.add\_embeddings}(project\_id, chroma\_entries)
\Else
    \State Optionally store embeddings only within MongoDB fields
\EndIf
\State \Return created\_chunks
\end{algorithmic}
\end{algorithm}
\FloatBarrier

\textbf{Implementation highlights:}
\begin{itemize}
    \item \textbf{Embedding model:} \texttt{sentence-transformers/all-MiniLM-L6-v2} (see embedding service and verification scripts).
    \item \textbf{Dual storage:} Metadata in MongoDB \texttt{text\_chunk} documents; vectors in ChromaDB via \texttt{backend/app/services/chroma\_service.py}.
    \item \textbf{Batch background tasks:} Celery jobs generate missing embeddings in \texttt{backend/app/tasks/indexing.py} (\texttt{batch\_generate\_embeddings\_task}).
\end{itemize}

\subsection{Pseudocode: Context-aware Retrieval (RAG)}

This step assembles a rich context from semantically similar chunks and optional entity/scene/relationship metadata.

\begin{algorithm}[H]
\caption{AssembleRAGContext}
\begin{algorithmic}[1]
\Require query \textit{Q}, \textit{project\_id}, optional \textit{file\_id}, max\_chunks \textit{k}
\Ensure \textit{context\_payload}: list of chunks, entities, scenes, relationships
\State query\_embedding $\leftarrow$ \texttt{EmbeddingService.generate\_single\_embedding}(Q)
\If{\texttt{ChromaService} available}
    \State vector\_hits $\leftarrow$ \texttt{ChromaService.search\_similar}(project\_id, query\_embedding, \texttt{n\_results} $=$ k)
\Else
    \State vector\_hits $\leftarrow$ \texttt{TextChunkRepository.vector\_search}(project\_id, query\_embedding, \texttt{limit} $=$ k)
\EndIf
\State top\_chunks $\leftarrow$ convert \textit{vector\_hits} to chunk payloads (\texttt{id}, \texttt{content}, \texttt{score}, \texttt{file\_id}, \texttt{positions})
\If{include\_entities}
    \State entities $\leftarrow$ \texttt{EntityRepository.get\_by\_project}(project\_id)
    \State entity\_matches $\leftarrow$ identify entities mentioned in \textit{top\_chunks} (\texttt{EntityExtractionService} or string matching)
\EndIf
\If{include\_scenes}
    \State scenes $\leftarrow$ \texttt{SceneRepository.find\_by\_project}(project\_id)
    \State mapped\_scenes $\leftarrow$ map chunks to scenes when positions overlap or proximity matches
\EndIf
\If{include\_relationships}
    \State relationships $\leftarrow$ \texttt{RelationshipRepository.find\_by\_project}(project\_id)
\EndIf
\State context\_payload $\leftarrow$ \{\texttt{"chunks"}: top\_chunks, \texttt{"entities"}: entity\_matches, \texttt{"scenes"}: mapped\_scenes, \texttt{"relationships"}: relationships (filtered)\}
\State \Return context\_payload
\end{algorithmic}
\end{algorithm}
\FloatBarrier

\textbf{Notes:}
\begin{itemize}
    \item \texttt{RAGContextService.assemble\_context} collects chunk-level results, enriches them with entity and scene metadata, and returns a compact context mapping suitable for prompt building (see \texttt{backend/app/services/rag\_context\_service.py}).
    \item \texttt{SearchService} implements hybrid logic and fallback ordering (Chroma $\rightarrow$ MongoDB \texttt{vectorSearch}).
\end{itemize}

\subsection{Pseudocode: AI Response Generation (RAG + LLM Call)}

This step orchestrates context assembly, prompt construction, LLM invocation, and persistence. It corresponds to \texttt{AIChatService.chat} and \texttt{CopilotService.generate\_suggestion}.

\begin{algorithm}[H]
\caption{GenerateAIResponseWithContext}
\begin{algorithmic}[1]
\Require \textit{user\_message}, \textit{project\_id}, optional \textit{file\_id}, \textit{user\_id}, \textit{max\_context\_chunks}
\Ensure \textit{assistant\_message}
\State \textbf{Step A: Assemble RAG context}
\State context $\leftarrow$ \texttt{RAGContextService.assemble\_context}(\,\texttt{query} $=$ user\_message, \texttt{project\_id} $=$ project\_id, \texttt{file\_id} $=$ file\_id, \texttt{max\_chunks} $=$ max\_context\_chunks\,)
\State \textbf{Step B: Format context into LLM prompt}
\State formatted\_context $\leftarrow$ \texttt{RAGContextService.format\_context\_for\_llm}(context)
\State system\_prompt $\leftarrow$ \texttt{AIChatService.SYSTEM\_PROMPT}
\State conversation\_history $\leftarrow$ \texttt{ConversationRepository.get\_latest\_messages}(conversation\_id or new)
\State llm\_messages $\leftarrow$ \texttt{build\_messages\_list}(system\_prompt, conversation\_history, user\_message, formatted\_context)
\Comment{system message, recent assistant/user turns, current user message, appended context}
\State \textbf{Step C: Call LLM}
\State assistant\_content $\leftarrow$ \texttt{GroqService.generate}(\texttt{messages} $=$ llm\_messages, \texttt{max\_tokens} $= \dots$, \texttt{model} $= \dots$)
\State \textbf{Step D: Postprocess and store}
\State cleaned $\leftarrow$ \texttt{clean\_text}(assistant\_content) \Comment{e.g., \texttt{CopilotService.\_clean\_suggestion}}
\State \texttt{MessageRepository.save}(role $=$ \texttt{assistant}, content $=$ cleaned, context\_used $=$ context)
\State update conversation metadata (message count, token usage, timestamps)
\State \Return cleaned
\end{algorithmic}
\end{algorithm}
\FloatBarrier

\textbf{Implementation notes:}
\begin{itemize}
    \item \texttt{AIChatService.chat} orchestrates assembling context, formatting prompts, invoking the LLM, and persisting messages.
    \item \texttt{CopilotService.generate\_suggestion} targets inline assistance: it infers writing style from RAG results, calls the LLM, then returns a cleaned suggestion and the context used (see \texttt{backend/app/services/copilot\_service.py}).
\end{itemize}

\subsection*{Complexity, Reliability, and Practical Concerns}
\label{subsec:complexity-reliability-practical}

\begin{itemize}
    \item \textbf{Chunking complexity:} $O(n)$ in the number of sentences per document. Overlap increases effective storage but preserves local context for downstream retrieval.
    \item \textbf{Embedding generation:} Dominated by model inference cost per chunk. Batch generation reduces overhead (indexing background tasks \texttt{batch\_generate\_missing\_embeddings}).
    \item \textbf{Vector search:} ChromaDB uses efficient ANN (HNSW) with sub-linear query cost; fallback MongoDB vector aggregation may be slower for large collections.
    \item \textbf{Fault tolerance:}
    \begin{itemize}
        \item Embedding/model calls may fail due to transient network or provider issues; services use \texttt{try/catch} with sensible fallbacks (see embedding service and verify scripts).
        \item RAG assembly respects available services and degrades gracefully (ChromaDB optional; code falls back to MongoDB vector search when unavailable).
    \end{itemize}
    \item \textbf{Reindexing \& incremental updates:} \texttt{TextExtractionService.reprocess\_file} deletes existing chunks and marks the file ready for reprocessing; batch background tasks handle any missing embeddings.
\end{itemize}

\begin{algorithm}[H]
\caption{ReprocessFileAndIncrementalUpdate}
\begin{algorithmic}[1]
\Require \textit{file\_id}, \textit{project\_id}
\Ensure File is reindexed and embeddings regenerated as needed
\State \texttt{TextChunkRepository.delete\_by\_file}(file\_id)
\State \texttt{FileRepository.mark\_ready\_for\_processing}(file\_id)
\State full\_text $\leftarrow$ \texttt{TextExtractionService.extract\_text\_from\_file}(file\_id)
\State chunks $\leftarrow$ \texttt{TextExtractionService.chunk\_text}(project\_id, file\_id, full\_text, \texttt{max\_chunk\_size}, \texttt{overlap})
\State \texttt{TextChunkRepository.bulk\_insert}(chunks)
\State enqueue \texttt{batch\_generate\_missing\_embeddings}(project\_id, file\_id)
\end{algorithmic}
\end{algorithm}
\FloatBarrier

\subsection*{Concrete File References (Key Implementation Points)}
\begin{itemize}
    \item \textbf{Text chunking and models}
    \begin{itemize}
        \item \texttt{backend/app/services/text\_extraction\_service.py} :: \texttt{chunk\_text()}, \texttt{\_create\_text\_chunk()}, \texttt{\_split\_into\_sentences()}
        \item \texttt{backend/app/repositories/text\_chunk\_repository.py} :: \texttt{create()}, \texttt{get\_by\_file()}, \texttt{get\_chunks\_without\_embeddings()}, \texttt{update\_embedding()}
    \end{itemize}
    \item \textbf{Embeddings and vector store}
    \begin{itemize}
        \item \texttt{backend/app/services/embedding\_service.py} :: \texttt{generate\_single\_embedding()}, \texttt{generate\_embeddings()}, \texttt{store\_text\_chunks\_with\_embeddings()}, \texttt{find\_similar\_chunks()}
        \item \texttt{backend/app/services/chroma\_service.py} :: \texttt{get\_or\_create\_collection()}, \texttt{add\_embeddings()}, \texttt{search\_similar()}
        \item \texttt{backend/app/tasks/indexing.py} :: \texttt{index\_document\_task()}, \texttt{batch\_generate\_embeddings\_task()}
    \end{itemize}
    \item \textbf{RAG and search}
    \begin{itemize}
        \item \texttt{backend/app/services/rag\_context\_service.py} :: \texttt{assemble\_context()}, \texttt{format\_context\_for\_llm()}
        \item \texttt{backend/app/services/search\_service.py} :: \texttt{semantic\_search()}, \texttt{find\_similar\_content()}, hybrid logic
    \end{itemize}
    \item \textbf{LLM orchestration and assistants}
    \begin{itemize}
        \item \texttt{backend/app/services/ai\_chat\_service.py} :: \texttt{chat()}
        \item \texttt{backend/app/services/copilot\_service.py} :: \texttt{generate\_suggestion()}, \texttt{\_gather\_story\_context()}
    \end{itemize}
\end{itemize}

% See Appendix E for call signatures
See Appendix~\ref{sec:appendix-e-pseudo-api} for pseudo-API call signatures.

\section{External APIs and SDKs}
\label{sec:external-apis}

The system integrates multiple third-party APIs and SDKs to support AI capabilities, authentication, and data persistence.

\begin{table}[H]
\centering
\begin{tabular}{|p{3cm}|p{4cm}|p{4cm}|p{3cm}|}
\hline
\textbf{API / SDK} & \textbf{Description} & \textbf{Purpose of Usage} & \textbf{Endpoint/ Function} \\
\hline
OpenAI API (v1) & Large language model services & Text generation and embeddings & /v1/chat/completions \\
\hline
FastAPI & Python web framework & REST API backend & @app.get(), @app.post() \\
\hline
PostgreSQL & Relational database system & User and project data storage & psycopg2 driver \\
\hline
VectorDB (FAISS/Pinecone) & Vector similarity engine & Semantic search and retrieval & similarity\_search() \\
\hline
\end{tabular}
\caption{External APIs and SDKs Used}
\end{table}

\section{Testing Details}
\label{sec:testing-details}

Testing played a critical role in validating the correctness, reliability, and stability of the implemented system. A combination of black box testing and unit testing strategies was adopted to ensure that the system meets its functional requirements.

\subsection{Black Box Testing}
\label{subsec:blackbox-testing}

Black box testing was used to validate the external behavior of the system without reference to internal implementation details. The focus was on verifying that system features produced correct outputs when interacting through public interfaces such as API endpoints.

\subsubsection{Purpose of Black Box Testing}

The objective of this testing was to ensure that system functionality aligned with the documented functional requirements by validating response codes, output structures, and observable system behavior.

\subsubsection{Testing Environment}

Testing was conducted using an isolated staging environment that included a running API server, a connected database, and a REST API testing client. All tests were executed externally without accessing source code.

% Insert environment screenshot

\subsubsection{Functional API Coverage}

The following functional modules were tested as black box components:

\begin{itemize}
    \item Authentication services (registration, login, token validation)
    \item Project management operations
    \item File handling operations
    \item Search and retrieval services
    \item Chat and conversational APIs
\end{itemize}

% Insert API success response screenshots

\subsubsection{Workflow and Scenario Testing}

End-to-end user workflows were validated through multi-step test scenarios, including project creation, file indexing, semantic search, and AI-assisted responses.

% Insert workflow execution screenshots

\subsubsection{Negative and Security Testing}

Invalid inputs, unauthorized requests, and forbidden access attempts were tested to verify that the system safely rejected improper usage without exposing internal errors.

% Insert 401 / 403 screenshots

\subsubsection{Error and Edge Case Validation}

Boundary conditions such as missing data, empty result sets, and invalid resource requests were evaluated to ensure stable system responses.

% Insert error response screenshots

\subsubsection{Black Box Test Evidence}

All executed test cases were recorded and maintained as structured documentation.
Figures~\ref{fig:bboxutest0} to~\ref{fig:bboxutest3} present the black box test evidence artifacts.

\begin{figure}[H]
    \centering
    \includegraphics[width=0.90\textwidth]{images/tests1.png}
    \caption{Black Box Test Evidence 1}
    \label{fig:bboxutest0}
\end{figure}

\begin{figure}[H]
    \centering
    \includegraphics[width=0.90\textwidth]{images/tests2.png}
    \caption{Black Box Test Evidence 2}
    \label{fig:bboxutest0b}
\end{figure}

\begin{figure}[H]
    \centering
    \includegraphics[width=0.90\textwidth]{images/bboxutest1.png}
    \caption{Black Box Test Evidence 3}
    \label{fig:bboxutest1}
\end{figure}

\begin{figure}[H]
    \centering
    \includegraphics[width=0.90\textwidth]{images/bboxutest2.png}
    \caption{Black Box Test Evidence 4}
    \label{fig:bboxutest2}
\end{figure}

\begin{figure}[H]
    \centering
    \includegraphics[width=0.90\textwidth]{images/bboxutest3.png}
    \caption{Black Box Test Evidence 5}
    \label{fig:bboxutest3}
\end{figure}

\subsubsection{Black Box Testing Summary}

The system demonstrated consistent and correct behavior for valid use cases and safely handled invalid or unauthorized operations.

\subsection{Unit Testing}
\label{subsec:unit-testing}

Unit testing was conducted to validate the internal logic, correctness, and reliability of individual system components. Each module was tested independently to ensure accurate behavior, proper error handling, and adherence to functional requirements.

\subsubsection{Implemented Unit Test Coverage}

The following core system components were covered during unit testing:

\begin{itemize}
    \item Authentication services (registration, login, token handling)
    \item File processing and metadata management
    \item Search and autocomplete operations
    \item Text extraction, chunking, and incremental processing
    \item Embedding generation, statistics, and similarity scoring
    \item LLM interaction services (GroqService)
    \item Edit proposal generation and diff processing
    \item Copilot assistance (suggestions, style analysis, prompt building)
    \item RAG-based context assembly
    \item Entity extraction and validation
    \item Relationship discovery and contextual analysis
\end{itemize}

\subsubsection{Testing Approach}

Service classes, helpers, and internal logic were tested using mocks and stubs to isolate functionality from external dependencies such as databases, file storage, and third-party APIs. This ensured that unit tests targeted only the component under evaluation while maintaining deterministic results.

\subsubsection{Failure and Boundary Validation}

Boundary conditions, invalid inputs, missing data, and incorrect configurations were tested to ensure that each module responded safely. Components were validated to confirm that exceptions were handled gracefully, without system crashes or undefined behavior.

% -------------------------------------------------------------
% ----------------------- TABLE 1 ------------------------------
% -------------------------------------------------------------

\subsubsection{Authentication and File Services Unit Tests}

\begin{figure}[H]
    \centering
    \includegraphics[width=0.90\textwidth]{images/wboxutest1.png}
    \caption{Unit Test Evidence 1}
    \label{fig:wboxutest1}
\end{figure}

\subsubsection{Search, Text Extraction, and Embedding Unit Tests}

\begin{figure}[H]
    \centering
    \includegraphics[width=0.90\textwidth]{images/wboxutest2.png}
    \caption{Unit Test Evidence 2}
    \label{fig:wboxutest2}
\end{figure}



\subsubsection{LLM, Editing, and Copilot Unit Tests}

\begin{figure}[H]
    \centering
    \includegraphics[width=0.90\textwidth]{images/wboxutest3.png}
    \caption{Unit Test Evidence 3}
    \label{fig:wboxutest3}
\end{figure}

\subsubsection{RAG Context, Entity Extraction, and Relationship Discovery Unit Tests}

\begin{figure}[H]
    \centering
    \includegraphics[width=0.90\textwidth]{images/wboxutest4.png}
    \caption{Unit Test Evidence 4}
    \label{fig:wboxutest4}
\end{figure}

\subsubsection{Unit Test Summary}

\begin{table}[h!]
\centering
\small
\begin{tabular}{|p{4cm}|p{3cm}|p{4cm}|}
\hline
\textbf{Category} & \textbf{Count} & \textbf{Notes} \\ \hline
Total Unit Tests Executed & 70 & Across all modules \\ \hline
Passing Tests & 64 &  \\ \hline
Failed Tests & 6 & Require fixes/mocking updates \\ \hline
Overall Pass Rate & 91\% & Stable with minor issues \\ \hline
\end{tabular}
\caption{Overall Unit Testing Summary}
\end{table}

Unit test execution evidence is shown in Figures~\ref{fig:wboxutest1} to~\ref{fig:wboxutest4}.

\subsubsection{Current Iteration Test Additions}

The current iteration introduced comprehensive test coverage for the three new features: Dialogue Generation, Character Repository, and Research Integration Agent. These tests validate both API-level integration and service-level unit functionality.

\textbf{Dialogue Generation Tests:}

\begin{itemize}
    \item \textbf{Integration Tests (5 test cases):}
    \begin{itemize}
        \item \texttt{test\_generate\_dialogue\_suggestion\_continue}: Validates continue-mode suggestion generation with character context
        \item \texttt{test\_generate\_dialogue\_suggestion\_rewrite}: Tests rewrite-mode suggestions for dialogue text
        \item \texttt{test\_generate\_suggestion\_validation\_error}: Verifies request validation for missing required fields
        \item \texttt{test\_generate\_suggestion\_unauthorized}: Ensures authentication is enforced
        \item \texttt{test\_track\_accept\_and\_reject\_suggestion}: Tests suggestion acceptance/rejection tracking
    \end{itemize}
    \item \textbf{Unit Tests (13 test cases):}
    \begin{itemize}
        \item \texttt{test\_generate\_suggestion\_success}: Validates end-to-end suggestion generation flow
        \item \texttt{test\_gather\_story\_context}: Tests character, location, and theme extraction from project context
        \item \texttt{test\_analyze\_writing\_style}: Validates POV, tense, and formality detection
        \item \texttt{test\_build\_prompt\_continue}: Tests prompt construction for continue mode with character details
        \item \texttt{test\_build\_prompt\_includes\_dialogue\_tone\_and\_character\_profile\_details}: Verifies character attribute injection into prompts
        \item \texttt{test\_analyze\_writing\_style\_detects\_informal\_tone\_from\_contractions}: Tests formality detection heuristics
        \item \texttt{test\_gather\_story\_context\_keeps\_only\_mentioned\_characters\_for\_dialogue}: Validates context filtering logic
        \item Additional tests for complete/rewrite modes, error handling, and sentence extraction
    \end{itemize}
\end{itemize}

\textbf{Character Repository Tests:}

\begin{itemize}
    \item \textbf{Integration Tests (3 test cases):}
    \begin{itemize}
        \item \texttt{test\_list\_project\_entities\_returns\_character\_profiles}: Validates entity listing endpoint returns character data
        \item \texttt{test\_list\_project\_entities\_character\_type\_filter}: Tests entity type filtering for character-only results
        \item \texttt{test\_list\_project\_entities\_unauthorized}: Ensures authentication is required
    \end{itemize}
    \item \textbf{Unit Tests (4 test cases):}
    \begin{itemize}
        \item \texttt{test\_entity\_model\_stores\_character\_attributes\_and\_role}: Validates character attribute persistence in entity model
        \item \texttt{test\_entity\_repository\_get\_by\_project\_and\_character\_type}: Tests repository filtering by project and entity type
        \item \texttt{test\_process\_uploaded\_file\_updates\_existing\_character\_metadata}: Validates incremental character profile updates
        \item \texttt{test\_process\_uploaded\_file\_creates\_new\_character\_when\_missing}: Tests new character creation flow
    \end{itemize}
\end{itemize}

\textbf{Research Integration Agent Tests:}

\begin{itemize}
    \item \textbf{Integration Tests (4 test cases):}
    \begin{itemize}
        \item \texttt{test\_research\_chat\_success\_returns\_response\_with\_sources}: Validates end-to-end research flow with source attribution
        \item \texttt{test\_research\_chat\_returns\_503\_when\_search\_unavailable}: Tests graceful degradation when Tavily API is unavailable
        \item \texttt{test\_research\_chat\_wraps\_unexpected\_errors\_as\_500}: Validates error handling for unexpected failures
        \item \texttt{test\_research\_chat\_unauthorized}: Ensures authentication is enforced
    \end{itemize}
    \item \textbf{Unit Tests (7 test cases):}
    \begin{itemize}
        \item \texttt{test\_search\_success\_maps\_results\_and\_clamps\_max\_results}: Validates search result mapping and result limiting
        \item \texttt{test\_search\_raises\_when\_api\_key\_missing}: Tests API key validation
        \item \texttt{test\_search\_reraises\_http\_status\_error}: Validates HTTP error propagation
        \item \texttt{test\_search\_reraises\_timeout\_exception}: Tests timeout handling
        \item \texttt{test\_format\_results\_for\_llm\_with\_results}: Validates LLM context formatting with source numbering
        \item \texttt{test\_format\_results\_for\_llm\_empty\_results}: Tests empty result handling
        \item \texttt{test\_is\_available}: Validates service availability checking
    \end{itemize}
\end{itemize}

\begin{table}[H]
\centering
\small
\begin{tabular}{|p{4.5cm}|p{2cm}|p{2cm}|p{2cm}|}
\hline
\textbf{Feature Module} & \textbf{Integration Tests} & \textbf{Unit Tests} & \textbf{Total} \\ \hline
Dialogue Generation & 5 & 13 & 18 \\ \hline
Character Repository & 3 & 4 & 7 \\ \hline
Research Integration Agent & 4 & 7 & 11 \\ \hline
\textbf{New Tests Total} & \textbf{12} & \textbf{24} & \textbf{36} \\ \hline
\end{tabular}
\caption{Current Iteration Test Coverage by Feature}
\end{table}

\subsection{Test Evidence}

Overall test execution output is shown in Figure~\ref{fig:test-exec-summary}.

\begin{figure}[H]
    \centering
    \includegraphics[width=0.90\textwidth]{images/tests-summary.png}
    \caption{Test Execution Results}
    \label{fig:test-exec-summary}
\end{figure}


\section{Test Summary}
\label{sec:test-summary}

The testing results showed that the system successfully handled valid requests and appropriately rejected invalid or unauthorized access attempts. The core workflows such as project creation, file uploading, semantic search, and AI-assisted responses were verified to function as expected during the current iteration.

All critical defects identified during testing were resolved, and remaining enhancements will be addressed in the next development phase.

\subsection{Overall Testing Statistics}

The current iteration expanded test coverage significantly with the addition of 36 new test cases for the three major features (Dialogue Generation, Character Repository, and Research Integration Agent). The comprehensive test suite now includes:

\begin{table}[H]
\centering
\begin{tabular}{|p{4cm}|p{3cm}|p{4cm}|}
\hline
\textbf{Test Category} & \textbf{Count} & \textbf{Coverage} \\ \hline
Total Backend Tests & 231 & Integration + Unit \\ \hline
Integration API Tests & ~70 & All API endpoints \\ \hline
Unit Service Tests & ~90 & Core services \\ \hline
Unit Repository Tests & ~30 & Data access layer \\ \hline
Unit Model Tests & ~41 & Data models \\ \hline
\textbf{Backend Pass Rate} & \textbf{91\%} & 210/231 passing \\ \hline
\textbf{Backend Code Coverage} & \textbf{32\%} & Statement coverage \\ \hline
\end{tabular}
\caption{Comprehensive Testing Statistics}
\end{table}

\textbf{Test Distribution by Module:}

\begin{itemize}
    \item \textbf{Authentication \& Authorization:} 6 tests (registration, login, token validation, access control)
    \item \textbf{Project Management:} 12 tests (CRUD operations, isolation, pagination)
    \item \textbf{File Operations:} 19 tests (upload, processing, content retrieval, reprocessing)
    \item \textbf{Search \& Retrieval:} 18 tests (semantic search, hybrid search, autocomplete, similarity)
    \item \textbf{Chat \& Conversation:} 13 tests (message handling, context assembly, conversation management)
    \item \textbf{Copilot \& Dialogue Generation:} 18 tests (suggestion generation, style analysis, context gathering)
    \item \textbf{Character Repository:} 7 tests (entity persistence, filtering, incremental updates)
    \item \textbf{Research Integration:} 11 tests (web search, synthesis, source attribution)
    \item \textbf{Entity Extraction:} 6 tests (character/location extraction, validation)
    \item \textbf{Relationship Discovery:} 9 tests (co-occurrence analysis, relationship classification)
    \item \textbf{RAG Context Assembly:} 5 tests (context retrieval, formatting, entity enrichment)
    \item \textbf{Embedding \& Vector Operations:} 3 tests (embedding generation, similarity calculation)
    \item \textbf{Text Processing:} 4 tests (extraction, chunking, incremental processing)
    \item \textbf{LLM Integration:} 12 tests (Groq service, HuggingFace service, genre-specific models)
    \item \textbf{Edit Proposals:} 12 tests (diff generation, proposal parsing, prompt building)
    \item \textbf{Indexing \& Background Tasks:} 8 tests (task management, status tracking, cancellation)
    \item \textbf{Position \& Scene APIs:} 29 tests (line retrieval, scene detection, entity mentions)
    \item \textbf{Relationships API:} 12 tests (discovery, network analysis, statistics)
    \item \textbf{Integration Workflows:} 3 tests (end-to-end project workflows, multi-file operations)
    \item \textbf{Health \& System:} 3 tests (health checks, API documentation availability)
\end{itemize}

\textbf{Code Coverage Analysis:}

The backend code coverage report indicates 32\% statement coverage across the codebase. Key coverage highlights:

\begin{itemize}
    \item \textbf{High Coverage Modules (>70\%):}
    \begin{itemize}
        \item Core configuration and models: 100\% (initialization modules)
        \item Research endpoint: 100\% (36/36 statements)
        \item Web search service: 92\% (47/51 statements)
        \item Copilot service: 87\% (130/149 statements)
        \item Edit proposal service: 85\% (70/82 statements)
        \item Groq service: 79\% (49/62 statements)
    \end{itemize}
    \item \textbf{Medium Coverage Modules (40-70\%):}
    \begin{itemize}
        \item AI chat service: 71\% (65/91 statements)
        \item RAG context service: 63\% (78/124 statements)
        \item Text extraction service: 56\% (140/248 statements)
        \item HuggingFace service: 51\% (35/69 statements)
    \end{itemize}
    \item \textbf{Lower Coverage Modules (<40\%):}
    \begin{itemize}
        \item Background indexing tasks: 0\% (302 statements untested - requires Celery infrastructure)
        \item Scene detection services: 0\% (268 statements - complex LLM-dependent logic)
        \item Search cache service: 0\% (210 statements - Redis-dependent caching layer)
        \item File service: 28\% (62/221 statements)
        \item Entity extraction service: 40\% (177/440 statements)
    \end{itemize}
\end{itemize}

\textbf{Testing Methodology:}

\begin{itemize}
    \item \textbf{Integration Tests:} Use pytest with async fixtures, test database isolation, and API client mocking
    \item \textbf{Unit Tests:} Employ comprehensive mocking (AsyncMock, Mock) to isolate service logic from external dependencies
    \item \textbf{Test Fixtures:} Shared fixtures for authentication, projects, files, and test data setup
    \item \textbf{Coverage Tools:} pytest-cov for statement coverage analysis, HTML reports for detailed line-by-line coverage
    \item \textbf{CI/CD Integration:} Tests run via pytest with coverage reporting; batch scripts provided for Windows/Linux environments
\end{itemize}

\subsection{Integration Testing}
\label{subsec:integration-testing}

\subsubsection{Summary}

Integration tests validated end-to-end behavior across API layers, databases, and services using the repository's pytest setup. The suite covers authentication, projects, files, search, chat, workflows, and health checks.
Figure~\ref{fig:apiintegtest1} provides integration API test evidence.

\begin{itemize}
    \item \textbf{Total integration tests:} ~70 (PHASE1 \& TESTING\_GUIDE) — reported as 100\% passing in the repo docs.
    \item \textbf{Pytest configuration:} Tests discovered under \texttt{backend/tests/}; notable files include \texttt{test\_auth\_api.py}, \texttt{test\_projects\_api.py}, \texttt{test\_files\_api.py}, \texttt{test\_search\_api.py}, \texttt{test\_chat\_api.py}, \texttt{test\_integration\_workflows.py}, and \texttt{test\_health.py}.
    \item \textbf{Notes:} Built from testing documentation (\texttt{WHITEBOX\_TESTING\_COMPLETE.md}, \texttt{PHASE1\_TESTING\_SUMMARY.md}, \texttt{TESTING\_GUIDE.md}, \texttt{TEST\_SUMMARY.md}) and concrete repo files (\texttt{conftest.py}, \texttt{pytest.ini}, run scripts). Exact counts may vary with updates.
\end{itemize}

\begin{figure}[H]
    \centering
    \includegraphics[width=0.90\textwidth]{images/apiintegtest1.png}
    \caption{Integration API Test Evidence}
    \label{fig:apiintegtest1}
\end{figure}

\subsubsection{Frontend Testing Infrastructure}

The frontend testing infrastructure is configured using Jest and React Testing Library for component and integration testing. The test configuration supports:

\begin{itemize}
    \item \textbf{Test Framework:} Jest with Next.js integration via \texttt{next/jest}
    \item \textbf{Test Environment:} jsdom for DOM simulation
    \item \textbf{Component Testing:} React Testing Library for component behavior validation
    \item \textbf{Coverage Collection:} Configured to collect coverage from all source files in \texttt{src/} directory
    \item \textbf{Module Resolution:} Path aliases configured (\texttt{@/} maps to \texttt{src/})
    \item \textbf{Transform Handling:} ESM package transformation for modern dependencies (geist, lucide-react, framer-motion, @radix-ui)
\end{itemize}

\textbf{Frontend Coverage Configuration:}

\begin{itemize}
    \item \textbf{Test Patterns:} \texttt{src/**/__tests__/**/*.{js,jsx,ts,tsx}} and \texttt{src/**/*.{test,spec}.{js,jsx,ts,tsx}}
    \item \textbf{Coverage Exclusions:} TypeScript definition files, layout components, global CSS
    \item \textbf{Coverage Output:} HTML reports generated in \texttt{frontend\_coverage/} directory
    \item \textbf{Setup Files:} \texttt{jest.setup.js} for test environment configuration
\end{itemize}

\textbf{Frontend Test Scope:}

The frontend testing infrastructure is prepared to validate:

\begin{itemize}
    \item Authentication flows (login, registration, password reset)
    \item Project management UI components
    \item File upload and management interfaces
    \item Chat and conversation UI components
    \item Copilot suggestion display and interaction
    \item Character repository visualization
    \item Research integration UI components
    \item Search and autocomplete functionality
    \item Error handling and loading states
    \item Responsive design and accessibility compliance
\end{itemize}

\textbf{Note:} Frontend test implementation is configured and ready for expansion in subsequent development phases. The current iteration focused on backend service validation and API integration testing to ensure core functionality stability before comprehensive UI testing.

\section{Current Iteration Features}
\label{sec:current-iteration-features}

This section describes three major features implemented in the current iteration: Dialogue Generation, Character Repository, and Research Integration Agent. These features extend the core RAG pipeline with specialized capabilities for creative writing assistance.

\subsection{Dialogue Generation Module}
\label{subsec:dialogue-generation}

The Dialogue Generation module enhances the copilot service with context-aware dialogue suggestions that maintain character voice consistency and narrative coherence. This module integrates character profiles, writing style analysis, and RAG-based context retrieval to generate authentic dialogue continuations.

\subsubsection{Module Summary}

The dialogue generation system operates as an extension of the \texttt{CopilotService} and provides:

\begin{itemize}
    \item \textbf{Character-aware suggestions:} Retrieves character profiles and attributes from the entity repository to ensure dialogue matches established character voices and personality traits.
    \item \textbf{Style-adaptive generation:} Analyzes the author's writing style (POV, tense, formality, sentence structure) and generates suggestions that match the existing narrative voice.
    \item \textbf{Context-aware prompting:} Uses RAG to retrieve similar text chunks and recent story events to inform dialogue generation with relevant narrative context.
    \item \textbf{Multi-mode support:} Supports continue, complete, and rewrite suggestion types for flexible writing assistance.
\end{itemize}

\textbf{Key implementation components:}
\begin{itemize}
    \item \texttt{backend/app/services/copilot\_service.py :: CopilotService.generate\_suggestion}
    \item \texttt{backend/app/services/copilot\_service.py :: CopilotService.\_gather\_story\_context}
    \item \texttt{backend/app/services/copilot\_service.py :: CopilotService.\_analyze\_writing\_style}
    \item \texttt{backend/app/services/copilot\_service.py :: CopilotService.\_build\_prompt}
    \item \texttt{backend/app/api/v1/endpoints/copilot.py :: generate\_suggestion}
\end{itemize}

\subsubsection{Algorithm: Context-Aware Dialogue Generation}

The dialogue generation algorithm assembles multi-layered context (characters, style, recent events) and constructs specialized prompts for the LLM to generate stylistically consistent dialogue.

\begin{algorithm}[H]
\caption{GenerateDialogueSuggestion}
\begin{algorithmic}[1]
\Require \textit{project\_id}, \textit{text\_before}, \textit{text\_after}, \textit{cursor\_position}, \textit{suggestion\_type}, \textit{max\_tokens}
\Ensure \textit{suggestion\_response}: generated text, context used, confidence score
\State \textbf{Step A: Extract immediate context}
\State immediate\_context $\leftarrow$ last 500 characters of \textit{text\_before}
\State \textbf{Step B: Gather story context}
\State entities $\leftarrow$ \texttt{EntityRepository.get\_by\_project}(project\_id)
\State mentioned\_characters $\leftarrow$ [\,];\; mentioned\_locations $\leftarrow$ [\,];\; mentioned\_themes $\leftarrow$ [\,]
\For{each entity in entities}
    \If{entity.name appears in immediate\_context (case-insensitive)}
        \If{entity.type == CHARACTER}
            \State append \{name, aliases, attributes\} to mentioned\_characters
        \ElsIf{entity.type == LOCATION}
            \State append \{name, attributes\} to mentioned\_locations
        \ElsIf{entity.type == THEME}
            \State append \{name\} to mentioned\_themes
        \EndIf
    \EndIf
\EndFor
\State \textbf{Step C: Retrieve style samples via RAG}
\State last\_sentences $\leftarrow$ extract last 1-2 sentences from immediate\_context
\If{length(last\_sentences) $>$ 20}
    \State query\_embedding $\leftarrow$ \texttt{EmbeddingService.generate\_single\_embedding}(last\_sentences)
    \State similar\_chunks $\leftarrow$ \texttt{EmbeddingService.semantic\_search}(project\_id, last\_sentences, limit $=$ 3, min\_similarity $=$ 0.3)
    \State style\_samples $\leftarrow$ [chunk.content[:200] for chunk in similar\_chunks]
\EndIf
\State \textbf{Step D: Analyze writing style}
\State sentences $\leftarrow$ split immediate\_context by '.'
\State avg\_sentence\_length $\leftarrow$ mean word count per sentence
\State pov $\leftarrow$ detect\_pov(immediate\_context) \Comment{first/second/third person}
\State tense $\leftarrow$ detect\_tense(immediate\_context) \Comment{past/present}
\State formality $\leftarrow$ detect\_formality(immediate\_context) \Comment{formal/informal based on contractions}
\State style\_characteristics $\leftarrow$ \{avg\_sentence\_length, pov, tense, formality\}
\State \textbf{Step E: Build specialized prompt}
\State story\_context $\leftarrow$ \{mentioned\_characters[:5], mentioned\_locations[0], mentioned\_themes[:3], style\_samples, style\_characteristics\}
\State prompt $\leftarrow$ \texttt{\_build\_prompt}(suggestion\_type, text\_before, text\_after, story\_context)
\Comment{Constructs prompt with character details, style guidelines, and examples}
\State \textbf{Step F: Generate suggestion via LLM}
\State raw\_suggestion $\leftarrow$ \texttt{GroqService.generate\_response}(prompt, temperature $=$ 0.7, max\_tokens)
\State \textbf{Step G: Clean and format output}
\State cleaned\_suggestion $\leftarrow$ remove artifacts (e.g., "Here's the continuation:", quotes, ellipsis)
\State \Return \{suggestion: cleaned\_suggestion, context\_used: story\_context, confidence: 0.8\}
\end{algorithmic}
\end{algorithm}
\FloatBarrier

\textbf{Implementation notes:}
\begin{itemize}
    \item The prompt construction varies by suggestion type (continue, complete, rewrite) and includes explicit instructions to match the author's voice, character consistency, and narrative style.
    \item Character attributes (personality traits, emotional tendencies, narrative role) are injected into the prompt to guide dialogue generation.
    \item Style analysis uses heuristics for POV detection (presence of "I", "you"), tense detection (past vs. present tense verbs), and formality (contraction usage).
    \item RAG-based style samples provide concrete examples of the author's writing patterns for the LLM to emulate.
\end{itemize}

\subsection{Character Repository Module}
\label{subsec:character-repository}

The Character Repository module provides persistent storage and retrieval of character profiles extracted from uploaded documents. Characters are stored as entities with rich attributes including personality traits, emotional tendencies, and narrative roles, enabling character-aware writing assistance.

\subsubsection{Module Summary}

The character repository system extends the entity extraction pipeline with character-specific functionality:

\begin{itemize}
    \item \textbf{Character profile persistence:} Stores character entities with structured attributes (personality traits, emotional tendencies, narrative role) in the entity model.
    \item \textbf{Project-scoped retrieval:} Provides API endpoints to list all characters within a project with optional type filtering.
    \item \textbf{Incremental updates:} Updates existing character profiles when new mentions are detected during file processing, merging mention counts and confidence scores.
    \item \textbf{Entity linking:} Links character mentions to text chunks and position indices for context-aware retrieval.
\end{itemize}

\textbf{Key implementation components:}
\begin{itemize}
    \item \texttt{backend/app/models/entity.py :: Entity} (with EntityType.CHARACTER)
    \item \texttt{backend/app/repositories/entity\_repository.py :: EntityRepository.get\_by\_project\_and\_type}
    \item \texttt{backend/app/api/v1/endpoints/genres.py :: list\_project\_entities} (with entity\_type filter)
    \item \texttt{backend/app/api/v1/endpoints/files.py :: process\_uploaded\_file} (character extraction and update logic)
    \item \texttt{backend/app/services/entity\_extraction\_service.py :: EntityExtractionService.extract\_entities\_from\_text}
\end{itemize}

\subsubsection{Algorithm: Character Profile Extraction and Update}

The character repository algorithm processes uploaded files, extracts character mentions, and either creates new character profiles or updates existing ones with incremental metadata.

\begin{algorithm}[H]
\caption{ExtractAndPersistCharacterProfiles}
\begin{algorithmic}[1]
\Require \textit{file\_id}, \textit{project\_id}
\Ensure Character entities created or updated in repository
\State \textbf{Step A: Extract text and create chunks}
\State file\_record $\leftarrow$ \texttt{FileRepository.get\_by\_id}(file\_id)
\State text\_content $\leftarrow$ file\_record.text\_content
\State chunks $\leftarrow$ \texttt{TextExtractionService.chunk\_text}(project\_id, file\_id, text\_content)
\State \textbf{Step B: Extract entities from text}
\State extracted\_entities $\leftarrow$ \texttt{EntityExtractionService.extract\_entities\_from\_text}(text\_content, project\_id, file\_id)
\Comment{Returns list of Entity objects with type, name, attributes, mention\_count, confidence\_score}
\State \textbf{Step C: Process each extracted character entity}
\For{each entity in extracted\_entities}
    \If{entity.type == CHARACTER}
        \State existing\_entity $\leftarrow$ \texttt{EntityRepository.get\_by\_name}(project\_id, entity.name)
        \If{existing\_entity exists}
            \State \textbf{Update existing character profile}
            \State new\_mention\_count $\leftarrow$ existing\_entity.mention\_count + entity.mention\_count
            \State new\_confidence $\leftarrow$ (existing\_entity.confidence\_score + entity.confidence\_score) / 2
            \State update\_payload $\leftarrow$ \{
                \Statex \hspace{4em} mention\_count: new\_mention\_count,
                \Statex \hspace{4em} confidence\_score: new\_confidence,
                \Statex \hspace{4em} last\_mentioned: entity.last\_mentioned,
                \Statex \hspace{4em} attributes: merge(existing\_entity.attributes, entity.attributes)
            \}
            \State \texttt{EntityRepository.update\_by\_id}(existing\_entity.id, update\_payload)
        \Else
            \State \textbf{Create new character profile}
            \State \texttt{EntityRepository.create}(entity)
        \EndIf
    \EndIf
\EndFor
\State \textbf{Step D: Store text chunks with embeddings}
\State \texttt{EmbeddingService.store\_text\_chunks\_with\_embeddings}(project\_id, file\_id, chunks)
\State \textbf{Step E: Discover relationships between characters}
\State \texttt{RelationshipDiscoveryService.discover\_relationships\_for\_project}(project\_id)
\State \textbf{Step F: Update file processing status}
\State \texttt{FileRepository.update\_processing\_status}(file\_id, status $=$ "completed")
\end{algorithmic}
\end{algorithm}
\FloatBarrier

\textbf{Implementation notes:}
\begin{itemize}
    \item Character attributes are stored in a flexible \texttt{attributes} dictionary field within the Entity model, supporting arbitrary key-value pairs (e.g., personality\_traits: ["witty", "guarded"], narrative\_role: "protagonist").
    \item The \texttt{get\_by\_project\_and\_type} repository method supports pagination (skip, limit) and type filtering for efficient character listing.
    \item Mention count and confidence score are incrementally updated when the same character is detected in multiple files or processing runs.
    \item Character profiles are linked to relationship entities for narrative graph construction.
\end{itemize}

\subsection{Research Integration Agent Module}
\label{subsec:research-integration}

The Research Integration Agent provides real-time web search capabilities integrated with AI-powered synthesis, enabling writers to research topics and receive contextually relevant, cited information without leaving the writing environment.

\subsubsection{Module Summary}

The research integration system combines web search with LLM-based synthesis:

\begin{itemize}
    \item \textbf{Real-time web search:} Integrates with Tavily Search API to retrieve current web information based on natural language queries.
    \item \textbf{AI-powered synthesis:} Uses Groq LLM to synthesize search results into coherent, writer-focused answers with inline source citations.
    \item \textbf{Source attribution:} Returns structured source metadata (title, URL, snippet) alongside AI-generated answers for verification and citation.
    \item \textbf{Conversation persistence:} Stores research queries and responses in conversation history for project context continuity.
\end{itemize}

\textbf{Key implementation components:}
\begin{itemize}
    \item \texttt{backend/app/services/web\_search\_service.py :: WebSearchService.search}
    \item \texttt{backend/app/services/web\_search\_service.py :: WebSearchService.format\_results\_for\_llm}
    \item \texttt{backend/app/api/v1/endpoints/research.py :: research\_chat}
    \item \texttt{backend/app/models/conversation.py :: ResearchRequest, ResearchResponse, SourceItem}
\end{itemize}

\subsubsection{Algorithm: Web Research and AI Synthesis}

The research integration algorithm orchestrates web search, result formatting, LLM synthesis, and conversation persistence to provide cited, contextual answers.

\begin{algorithm}[H]
\caption{PerformResearchAndSynthesize}
\begin{algorithmic}[1]
\Require \textit{user\_query}, \textit{project\_id}, \textit{user\_id}, \textit{max\_sources}
\Ensure \textit{research\_response}: AI answer, source list, conversation metadata
\State \textbf{Step A: Validate web search availability}
\If{\texttt{WebSearchService.is\_available()} == False}
    \State \Return HTTPException(503, "Web search service not configured")
\EndIf
\State \textbf{Step B: Execute web search}
\State search\_results $\leftarrow$ \texttt{WebSearchService.search}(
    \Statex \hspace{4em} query $=$ user\_query,
    \Statex \hspace{4em} max\_results $=$ max\_sources,
    \Statex \hspace{4em} search\_depth $=$ "basic")
\Comment{Returns list of SearchResult objects with title, url, snippet, score}
\State \textbf{Step C: Format search results as LLM context}
\State search\_context $\leftarrow$ \texttt{WebSearchService.format\_results\_for\_llm}(search\_results)
\Comment{Formats as "[Source 1] Title\textbackslash nURL: ...\textbackslash nSnippet\textbackslash n"}
\State \textbf{Step D: Build research system prompt}
\State system\_prompt $\leftarrow$ "You are an intelligent research assistant for writers. Synthesize information from search results, cite sources as [Source N], focus on actionable information."
\State \textbf{Step E: Construct LLM messages}
\State llm\_messages $\leftarrow$ [
    \Statex \hspace{4em} \{role: "system", content: system\_prompt\},
    \Statex \hspace{4em} \{role: "user", content: search\_context + "\textbackslash n\textbackslash nUser Question: " + user\_query + "\textbackslash n\textbackslash nPlease provide a comprehensive answer based on the search results above."\}]
\State \textbf{Step F: Generate AI synthesis}
\State ai\_response $\leftarrow$ \texttt{GroqService.continue\_conversation}(
    \Statex \hspace{4em} messages $=$ llm\_messages,
    \Statex \hspace{4em} temperature $=$ 0.5,
    \Statex \hspace{4em} max\_tokens $=$ 2500)
\State \textbf{Step G: Persist conversation}
\State conversation $\leftarrow$ \texttt{AIChatService.\_create\_conversation}(user\_id, project\_id)
\State \texttt{AIChatService.\_store\_message}(conversation.id, role $=$ USER, content $=$ user\_query)
\State assistant\_msg $\leftarrow$ \texttt{AIChatService.\_store\_message}(
    \Statex \hspace{4em} conversation.id,
    \Statex \hspace{4em} role $=$ ASSISTANT,
    \Statex \hspace{4em} content $=$ ai\_response,
    \Statex \hspace{4em} context\_used $=$ \{type: "web\_research", sources\_count: length(search\_results)\},
    \Statex \hspace{4em} model $=$ "llama-3.3-70b-versatile")
\State \texttt{ConversationRepository.increment\_message\_count}(conversation.id)
\State \textbf{Step H: Build source items}
\State sources $\leftarrow$ [SourceItem(title $=$ r.title, url $=$ r.url, snippet $=$ r.snippet[:200]) for r in search\_results]
\State \textbf{Step I: Return research response}
\State \Return ResearchResponse(
    \Statex \hspace{4em} conversation\_id $=$ conversation.id,
    \Statex \hspace{4em} message $=$ MessageResponse(assistant\_msg),
    \Statex \hspace{4em} sources $=$ sources,
    \Statex \hspace{4em} search\_query $=$ user\_query)
\end{algorithmic}
\end{algorithm}
\FloatBarrier

\textbf{Implementation notes:}
\begin{itemize}
    \item The Tavily Search API is configured via \texttt{TAVILY\_API\_KEY} environment variable; the service gracefully degrades if unavailable.
    \item Search results are limited to 1-10 sources (configurable via \texttt{max\_sources} parameter) to balance context richness and LLM token limits.
    \item The research system prompt explicitly instructs the LLM to cite sources using [Source N] notation and avoid fabricating information.
    \item Conversation persistence enables research history tracking and potential future features like research session resumption.
    \item Source snippets are truncated to 200 characters in the response payload for UI display efficiency.
\end{itemize}

